\chapter{PROPOSED SYSTEM}

\paragraph\
JASIPA is a web application in which a set of cooperating agents prepares the evidence for a hiring decision and an HR reviewer makes the decision. This chapter describes the system, its users, its advantages and its boundary.

\section{Description of Proposed System}

\paragraph\
Figure \ref{Fwk1} shows the architecture. The React frontend has an HR workspace and a candidate portal. Both call a FastAPI backend that enforces authentication and role-based access. Behind the API sit the LangGraph orchestrator and its agents, the QuestionGenerationAgent, the code sandbox, the resume parser and the bias checker. All state lives in a relational database (SQLite or PostgreSQL), uploaded resume files live in private storage, and an MCP server exposes ATS records to tools.

\begin{figure}[H]
\centering
\resizebox{0.85\textwidth}{!}{\begin{tikzpicture}[node distance=5mm and 7mm]
\node[title] (t1) {Frontend (React 19 + Vite)};
\node[box, below=2mm of t1, text width=26mm] (hr) {HR workspace\\\scriptsize Dashboard, Candidates,\\ Jobs, Bias, Audit};
\node[box, right=of hr, text width=26mm] (cp) {Candidate portal\\\scriptsize Resume, Roles,\\ Assessments};
\node[group, fit=(t1)(hr)(cp)] (fe) {};
\node[title, below=14mm of fe.south] (t2) {Backend API (FastAPI, JWT, RBAC)};
\node[box, below=2mm of t2, text width=18mm] (r1) {Auth \&\\ Candidates};
\node[box, right=3mm of r1, text width=18mm] (r2) {Jobs \&\\ Assessments};
\node[box, right=3mm of r2, text width=18mm] (r3) {Pipeline \&\\ Reviews};
\node[box, left=3mm of r1, text width=18mm] (r0) {Bias, Audit,\\ Dashboard};
\node[group, fit=(t2)(r0)(r3)] (api) {};
\node[agent, below=16mm of r1, text width=26mm] (lg) {LangGraph\\ orchestrator};
\node[agent, right=6mm of lg, text width=30mm] (qg) {QuestionGeneration-\\Agent};
\node[agent, left=6mm of lg, text width=26mm] (bias) {Bias checker\\\scriptsize Fisher, Wilson, 4/5};
\node[box, below=7mm of qg, text width=30mm] (sb) {Code sandbox\\\scriptsize isolated interpreter};
\node[box, below=7mm of lg, text width=26mm] (parse) {Resume parser\\\scriptsize PDF/DOCX/Kaggle HTML};
\node[box, below=7mm of bias, text width=26mm] (llm) {Gemini (optional)\\\scriptsize deterministic fallback};
\node[store, below=22mm of parse, text width=34mm] (db) {SQLite / PostgreSQL\\\scriptsize ATS records, audit ledger};
\node[store, left=8mm of db, text width=24mm] (csv) {Kaggle Resume\\ Dataset (CSV)};
\node[store, right=8mm of db, text width=24mm] (files) {Resume files\\\scriptsize private storage};
\node[box, right=8mm of files, text width=20mm] (mcp) {MCP server\\\scriptsize stdio + REST};
\draw[arr] (fe) -- node[lbl, right] {HTTPS / JSON} (api);
\draw[arr] (r3.south) -- (lg);
\draw[arr] (r2.south) -- (qg);
\draw[arr] (r0.south) -- (bias);
\draw[arr] (qg) -- (sb);
\draw[arr] (lg) -- (parse);
\draw[arr] (lg) -- (llm);
\draw[arr] (lg) -- (bias);
\draw[arr] (parse) -- (db);
\draw[arr] (csv) -- node[lbl, above] {import} (db);
\draw[arr] (parse) -- (files);
\draw[arr] (mcp.south) -- ++(0,-5mm) -| node[lbl, below, pos=0.25] {ATS tools} (db.south);
\end{tikzpicture}
}
\caption{System architecture}
\label{Fwk1}
\end{figure}

\paragraph\
The main modules are:
\begin{itemize}
\item \textbf{Resume ingestion:} validates and parses PDF/DOCX uploads and Kaggle dataset HTML into section-structured text and extracted fields.
\item \textbf{Rubric and question generation:} creates a versioned, JD-anchored rubric for every job and an assessment for every rubric version.
\item \textbf{Evaluation swarm:} the Resume-Screener, Skill-Assessor and Culture-Fit agents, merged by the Panel Coordinator.
\item \textbf{Bias checker:} cohort statistics bound to each panel decision.
\item \textbf{Decision and audit:} human decisions, two-reviewer escalation and the hash-chained audit ledger.
\item \textbf{MCP connector:} ATS tools over stdio JSON-RPC and REST.
\end{itemize}

\section{Description of Target Users}
\paragraph\
\textbf{HR reviewers and recruiters} create jobs, adjust rubrics, add candidates, run evaluations, review evidence and record decisions. \textbf{Candidates} register, upload resumes, apply and take assessments; they see their progress but never scores, answer keys or other candidates' data. \textbf{Auditors and compliance staff} use the audit log and bias monitor. The design principles followed are separation of duties (agents advise, people decide), least privilege (role-scoped APIs, sandboxed code), traceability (versioned rubrics, hash-chained logs) and graceful degradation (deterministic fallbacks for every LLM call).

\section{Advantages/Applications of Proposed System}
\paragraph\
\begin{itemize}
\item \textbf{Consistent, explainable scoring:} every applicant to a rubric version answers the same questions and is scored against the same weights, with evidence lines for each score.
\item \textbf{Measured skill:} code is executed against hidden tests, so a plausible-looking but wrong solution does not score well.
\item \textbf{Honest fairness reporting:} small cohorts are reported as insufficient data rather than as fair, and a disparity is flagged only when a defined screening criterion is triggered.
\item \textbf{Human accountability:} no automated rejection; escalated cases need two reviewers; every decision is bound to the evidence and bias check it was based on.
\item \textbf{Works offline:} without an LLM key the system runs fully on deterministic strategies.
\end{itemize}
Applications include campus and lateral hiring for technical and non-technical roles, and auditing an organisation's own screening process.

\section{Scope (Boundary of proposed system)}
\paragraph\
\textbf{In scope (implemented):} resume ingestion, rubric versioning, assessment generation and regeneration, sandboxed Python coding questions, the three-agent swarm with a Panel Coordinator, per-decision bias checks, human decisions with two-reviewer escalation, the audit ledger, the MCP connector, Kaggle dataset import, hiring-outcome recording and descriptive analytics.

\textbf{Out of scope or not implemented:} interview scheduling, video interviews, OCR of scanned resumes, coding languages other than Python, single sign-on and user management for HR accounts, and any automatic learning from decisions or outcomes. The system does not hire or reject anyone.
